\section{Реализация ядра физического движка}

\subsection{Структура модулей}

Исходный код проекта разбит на четыре модуля (рис.~\ref{fig:module_arch}):

\begin{figure}[h]
\centering
\begin{tikzpicture}[
  font=\small,
  mod/.style={draw, rounded corners=4pt, minimum width=4cm,
              minimum height=0.9cm, align=center, fill=white},
  arr/.style={-{Stealth[length=5pt]}, thick},
]
\node[mod, fill=blue!12] (comp) at (0,0)  {\texttt{compute/}: GPU-абстракция};
\node[mod, fill=green!12](core) at (0,1.4){\texttt{core/}: физика};
\node[mod, fill=yellow!12](scene) at (-2.5,2.8){\texttt{scene/}: загрузка сцен};
\node[mod, fill=orange!12](bind) at (2.5,2.8) {\texttt{bindings/}: Python};

\draw[arr] (core.south)  -- (comp.north);
\draw[arr] (scene.south) -- (core.north);
\draw[arr] (bind.south)  -- (core.north);
\end{tikzpicture}
\caption{Зависимости между модулями проекта}
\label{fig:module_arch}
\end{figure}

Модуль \texttt{core/} реализует всю физическую логику и зависит только от \texttt{compute/}.
Внешние зависимости (Eigen, spdlog, pugixml) не проникают в GPU-ядра:
lint-проверка при сборке запрещает их использование в коде устройства.

\subsection{Интегратор}

Интегратор реализует симплектический метод Эйлера с дроблением шага
на $n_{\text{sub}}$ подшагов.
Каждый подшаг запускает одно GPU-ядро, обрабатывающее все тела параллельно.

Симплектическая схема обновляет скорость \textit{до} позиции,
что гарантирует сохранение энергии системы при отсутствии диссипации:
\begin{align}
  \mathbf{v}(t + h) &= \mathbf{v}(t) + h\,\mathbf{g}, \\
  \mathbf{p}(t + h) &= \mathbf{p}(t) + h\,\mathbf{v}(t + h).
\end{align}
Кватернион ориентации обновляется через произведение с кватернионом
угловой скорости и нормализуется после каждого шага.

Тела с нулевой обратной массой ($m^{-1} = 0$) и тела с установленным
признаком \textit{Sleeping} пропускаются в ядре без ветвления по всему массиву.

\subsection{Широкая фаза: LBVH на GPU}
\label{sec:broadphase}

Широкая фаза реализует параллельную линейно-кодируемую BVH
(Linearly-encoded Bounding Volume Hierarchy, LBVH)~\cite{karras2012lbvh}.

Алгоритм включает четыре последовательных GPU-прохода (рис.~\ref{fig:lbvh_pipeline}):

\begin{figure}[h]
\centering
\begin{tikzpicture}[
  font=\small,
  pstep/.style={draw, rounded corners=3pt, fill=blue!10,
               minimum width=6cm, minimum height=0.7cm, align=center},
  arr/.style={-{Stealth[length=5pt]}, thick},
]
\node[pstep] (morton) at (0,0)    {1. Вычисление кодов Мортона по позициям тел};
\node[pstep] (sort)   at (0,-1.1) {2. Сортировка по ключу (Morton + индекс тела)};
\node[pstep] (tree)   at (0,-2.2) {3. Построение дерева (алгоритм Карраса)};
\node[pstep] (refit)  at (0,-3.3) {4. Обновление AABB снизу вверх (атомики)};

\draw[arr] (morton.south) -- (sort.north);
\draw[arr] (sort.south)   -- (tree.north);
\draw[arr] (tree.south)   -- (refit.north);
\end{tikzpicture}
\caption{Этапы построения LBVH на GPU}
\label{fig:lbvh_pipeline}
\end{figure}

\textbf{Коды Мортона.}
Каждая позиция тела нормируется относительно AABB сцены и квантуется
в три 10-битных координаты; их интерливинг образует 30-битный код Мортона.
Код Мортона кодирует пространственную близость: телá с близкими кодами
расположены в одном регионе пространства.
В качестве ключа сортировки используется 64-битное слово
$(\text{Morton} \ll 32) \,|\, i$, что делает сортировку стабильной
при совпадении кодов.

\textbf{Построение дерева.}
Алгоритм Карраса строит бинарное дерево из $n$ листьев и $n-1$ внутренних узлов
за один проход: каждый внутренний узел $i$ определяет свой диапазон листьев
через наибольший общий префикс кодов Мортона.
Вся структура дерева вычисляется параллельно~--- по одному потоку на узел.

\textbf{Обновление AABB.}
AABB внутренних узлов вычисляются снизу вверх.
Каждый поток, завершив обработку листа, атомарно увеличивает счётчик родительского узла;
второй поток, достигший родителя, начинает его обработку.
AABB узлов хранятся в формате fp16 для снижения потребления памяти вдвое.

После построения дерева каждый лист обходит дерево, эмитируя кандидатные пары $(i, j)$,
чьи AABB перекрываются.

\subsection{Узкая фаза}

Узкая фаза обрабатывает кандидатные пары параллельно.
Для каждой пары по типам форм выбирается один из алгоритмов (таблица~\ref{tab:narrowphase}).

\begin{table}[h]
\centering
\caption{Алгоритмы узкой фазы}
\label{tab:narrowphase}
\begin{tabular}{ll}
\toprule
\textbf{Пара форм} & \textbf{Алгоритм} \\
\midrule
Сфера--Сфера              & Аналитический (расстояние между центрами) \\
Сфера--Параллелепипед     & Ближайшая точка на OBB \\
Параллелепипед--Параллелепипед & SAT (15 осей разделения) + извлечение многоточечного контакта \\
Выпуклый многогранник~--~* & GJK (расстояние) + EPA (глубина проникновения) \\
*~--~Треугольная сетка    & Обход BVH + тест примитив--треугольник \\
\bottomrule
\end{tabular}
\end{table}

GJK и EPA реализованы с фиксированными массивами фиксированного размера
(64 вершины, 128 граней, 32 итерации для EPA),
чтобы исключить динамическое выделение памяти в GPU-ядре.

\subsection{XPBD-решатель}
\label{sec:xpbd_impl}

Решатель реализует последовательный Гаусс--Зейдель (Sequential Gauss--Seidel) по уравнениям
раздела~2.4.
Алгоритм~\ref{alg:xpbd} описывает один шаг решателя для набора контактов и сочленений.

\begin{lstlisting}[style=pseudocode,
  caption={Один шаг XPBD-решателя},
  label={alg:xpbd}]
Вход: контакты C, сочленения J, тела B, шаг h, число итераций $n$
  обнулить $\lambda_c \leftarrow 0$ для каждого контакта
  для $k = 1 \ldots n$:
    для каждого контакта $c \in C$ с глубиной $d$, нормалью $\hat{n}$:
      $w \leftarrow w_A(\hat{n}) + w_B(\hat{n})$   // обобщ. обратная масса
      $\Delta\lambda \leftarrow -(d + \tilde{\alpha}\,\lambda_c)\;/\;(w + \tilde{\alpha})$
      $\lambda_c \leftarrow \max(\lambda_c + \Delta\lambda,\;0)$  // $\lambda \geq 0$
      скорректировать позицию и ориентацию тел A и B
    для каждого сочленения $j \in J$:
      применить позиционные и угловые ограничения типа $j$
      применить PD-мотор и упоры (если задан тип Revolute/Prismatic)
  обновить скорости: $\mathbf{v} \leftarrow (\mathbf{x} - \mathbf{x}_{\mathrm{prev}})/h$
\end{lstlisting}

Для сочленений типа \textit{Fixed} и \textit{Ball} применяются
позиционные ограничения (шаги~1--4 без проекции).
Для \textit{Revolute} и \textit{Prismatic} добавляются угловые ограничения,
а также проверка упоров и поправка от PD-мотора.

Квотернионы нормализуются после каждой итерации~\eqref{eq:quat-update}
для предотвращения численного дрейфа.

По завершении всех итераций скорости тел обновляются
в соответствии с уравнением~\eqref{eq:vel-update}.

\subsection{Сцена и конвейер загрузки активов}
\label{sec:scene}

Модуль \texttt{scene/} реализует загрузку сцен из форматов SDF/URDF.
SDF (Simulation Description Format) является родным форматом симулятора Gazebo
и поддерживает вложенные модели, сочленения, материалы и коллизионные меши.

Конвейер загрузки включает следующие этапы:
\begin{enumerate}
  \item Разбор XML через библиотеку \textbf{libsdformat};
    URDF конвертируется в SDF внутри библиотеки.
  \item Для невыпуклых коллизионных мешей выполняется
    \textbf{выпуклая декомпозиция} инструментом V-HACD~4.1:
    меш разбивается на $k \leq k_{\max}$ выпуклых оболочек,
    каждая из которых регистрируется как отдельная форма.
  \item Треугольные меши статических тел загружаются как \textbf{BVH-деревья}
    (CPU-построение по медианному разбиению, однократная загрузка на GPU).
\end{enumerate}

Координатная система SDF~--- $z$ вверх; гравитация задаётся как $(0, 0, -9{,}81)$\,м/с²
без дополнительного преобразования осей.

\subsection{Выводы по разделу}

Ядро физического движка опирается на три архитектурных решения:
LBVH обеспечивает $O(n \log n)$ широкую фазу, полностью параллельную на GPU;
SoA-формат данных и слой \texttt{compute/} отделяют физическую логику от конкретного GPU API;
XPBD-решатель обеспечивает единообразную обработку контактов и сочленений
одним итерационным алгоритмом.
